% Track r4 -- the bridge T = nJ(n) between the two readings.
% Verification: investigations/two_readings/verify/r4_bridge.py (claims r4-1 ... r4-24).
% New references: sections/bib_r4.tex.
\section{\texorpdfstring{The bridge $T=nJ(n)$ between the two readings}{The bridge T = nJ(n) between the two readings}}
\label{r4:sec:main}

Two sections of this note meet in one element of $\Cl(3,3)$. The clock $n\in V^+$ of \cref{r2:sec:main} determines the complex structure
$J(n)=Tn$ of the Dirac reading (\cref{r2:prop:J,r2:thm:i}), while the octant representation of the field in \cref{r3:sec:main} uses $T=e_1e_2e_3$
itself, through $-T$, as the duality operator (\cref{r3:prop:parity,r3:prop:odd}). Since $n$ commutes with $T$ and $n^2=\id$,
$T=nJ(n)=J(n)n$ for every clock. This section asks whether this identity can be built into both readings at once while their distinct
Lorentz representations are kept apart, also when the clock varies or the temporal axes are permuted, and it answers the two points left open
in \cref{r3:sec:main}: the remark on the spatial derivatives of the octant Maxwell operator after \cref{r3:thm:maxwell}, and the paragraph ``The
main obstruction'' of \cref{r3:sec:answer}. In brief:
(1) $T$ is invariant and $J(n)$ covariant under every rotation of the temporal triple and under the diagonal rotations; an odd permutation of the
temporal axes reverses $T$ and $J(n)$ together, and with them the mass form and the duality operator, and turns the octant map into its
textbook-orientation form; reversing the clock is a half-turn, which fixes $T$, and the charge conjugation of \cite{ES17,RE} is that half-turn
composed with $T$ (\cref{r4:sec:identity});
(2) $T$ is at once the Dirac mass form, the Dirac parity and the lift of the frame parity, and the duality on the octants; $J(n)$ acts on the
octants as a rotation of the transverse field about $n$ and maps the longitudinal field to the two poles that carry the Gauss laws
(\cref{r4:sec:roles});
(3) the spin lift of $\lint(n)$ and the octant representation $\rho_{\rm oct}$ are inequivalent on the same module, their joint commutant is $\R\id$,
no covariant bilinear coupling exists, and the covariant trilinear couplings of a Grassmann Dirac field to the field octants form a
two-dimensional space, the Pauli term and its dual, written through $T$ and $J(n)$ (\cref{r4:sec:reps});
(4) the two first-order operators use conjugate Pauli triples, related only by elements of mixed parity, and neither derives spatial
differentiation from temporal directions (\cref{r4:sec:operators});
(5) the field reading does not depend on the clock and the particle reading does, and $T$ is, up to scale, the only element that is the clock
times a Dirac complex structure for every clock (\cref{r4:sec:clock}). \Cref{r4:sec:answer} collects the verdicts.

\paragraph{Conventions and units.} Those of \cref{r2:sec:main,r3:sec:main}: $V=\R^{3,3}$, $\eta=\operatorname{diag}(1,1,1,-1,-1,-1)$, the
realisation of eq.~(1) of \cite{ES17}, $s_a:=e_{a+3}$, $T:=e_1e_2e_3$, $\Gamma:=\omega=e_1\cdots e_6$, the binding root $Ge_i=e_{i+3}$,
$Ge_{i+3}=-e_i$, the diagonal rotations $L_i=\tfrac12(e_{j+3}e_{k+3}-e_je_k)$ and boosts $K_i=\tfrac12(e_ke_{j+3}-e_je_{k+3})$ of $\lG$ for cyclic
$(i,j,k)$, acting on $V$ by the commutator, and the octant notation $|\mathrm{vac}\rangle$, $|\mathrm{full}\rangle$, $f_k$, $g_m$, $\Phi(X)=\bE\cdot f+\bB\cdot g$ of
\cref{r3:prop:octants,r3:lem:rot,r3:thm:phi}; for $v\in\R^3$ put $f(v)=\sum_kv_kf_k$ and $g(v)=\sum_mv_mg_m$. The field point is
$X=\sum_i(B_ie_i+E_ie_{i+3})$ with $\bB$ minus the textbook magnetic field, $\Btext=-\bB$. A \emph{clock} is a unit $n\in V^+$ with an oriented
orthonormal frame $(n,u_1,u_2)$ of $V^+$, so $nu_1u_2=T$ (\cref{r2:prop:J}); $\gamma^\mu_n=(n,s_1,s_2,s_3)$, $\lint(n)$ is spanned by the
bivectors of $E_n=\R n\oplus V^-$, $\Omega_n:=ns_1s_2s_3$, and $\phi_n:\lG\to\lint(n)$ is the dictionary of \cref{r2:prop:dict},
$L_i\mapsto s_j\wedge s_k$, $K_i\mapsto s_i\wedge n$. A bivector $a\wedge b$ of orthogonal vectors is the Clifford element $\tfrac12ab$, and its spin
lift is left multiplication by that element, so the Dirac reading's Lorentz representation of $\lG$ on $S$ is
\[
\rho_n:=\text{spin}\circ\phi_n,\qquad \rho_n(L_i)=\tfrac12s_js_k,\qquad \rho_n(K_i)=\tfrac12s_in ,
\]
and the field reading's is $\rho_{\rm oct}(L_i)=L_i$, $\rho_{\rm oct}(K_i)=Q_i=-TL_i$ (\cref{r3:prop:odd}). Further $t_i:=\tfrac12e_je_k$ spans the
rotations $\so3_+$ of the temporal triple alone, and $A_i:=\tfrac12(e_je_k+e_{j+3}e_{k+3})$ rotates the two triples in opposite senses
(\cref{r3:lem:rot}); $L(n):=\sum_in_iL_i$, $A(n):=\sum_in_iA_i$. Dimensions: $n$, $J(n)$, $T$, $\Gamma$, all Clifford elements and all
intertwiners are dimensionless; $[x^a]=[x^0]=[ct]=\mathrm m$ and $[\partial_\mu]=\mathrm m^{-1}$; in SI units a Dirac field has $[\psi]=\mathrm m^{-3/2}$ and
a Lagrangian density $\mathrm{J\,m^{-3}}$ (\cref{r1:eq:LD}); the field enters as $X_{\rm EM}=(c\bB;\bE)$ with $[E_i]=[cB_i]=\mathrm{V\,m^{-1}}$; a
magnetic moment $\mu$ has $[\mu]=\mathrm{J\,T^{-1}}$ and an electric dipole moment $d$ has $[d]=\mathrm{C\,m}$. Computations use Heaviside--Lorentz units
with $\hbar=c=1$. Identifiers r4-$n$ refer to \texttt{verify/r4\_bridge.py}, which checks every finite-dimensional and symbolic claim below
in exact arithmetic (r4-1 re-verifies the conventions).

\subsection{The identity and its covariance}
\label{r4:sec:identity}

\begin{proposition}[The identity]
\label{r4:prop:identity}
For every $n\in V^+$: $n^2=|n|^2\id$, $J(n):=Tn=nT=n_1e_2e_3+n_2e_3e_1+n_3e_1e_2$, $J(n)^2=-|n|^2\id$, and $nJ(n)=J(n)n=|n|^2\,T$. Hence for every clock
\[
T=nJ(n)=J(n)n,\qquad J(n)=Tn=nT,\qquad n^2=\id,\qquad J(n)^2=T^2=-\id .
\]
$J(n)$ commutes with $n,s_1,s_2,s_3$, and $n\mapsto J(n)$ is linear and injective.
\end{proposition}

\begin{proof}
By the blade rule (Lemma~2.1 of \cite{ES17}) $T$ commutes with each $e_i$ and anticommutes with each $s_a$. So $J(n)^2=TnTn=T^2n^2$ and
$nJ(n)=nnT=|n|^2T$; $T$ and $n$ each anticommute with $s_a$, so $J(n)$ commutes with $s_a$ (r4-2, symbolic in $n$ and on the rational parametrisation
$n=(2a,2b,1-a^2-b^2)/(1+a^2+b^2)$ of the clock sphere).
\end{proof}

Restated: \cref{r2:prop:J} and Remark~4.2 of \cite{ES17} ($J=\gamma^0T$), re-verified (r4-2). Proposition~5.1 of \cite{FMB} is confirmed.

\begin{proposition}[Rotations]
\label{r4:prop:rot}
\begin{enumerate}[label=(\roman*),leftmargin=2.2em]
\item For $Y\in\so3_+\oplus\so3_\Delta$, $[Y,T]=0$ and $[Y,J(n)]=J([Y,n])$, with $[Y,n]\in V^+$. Hence for the spin lift $R$ of a rotation of the temporal
triple alone, or of a diagonal rotation, $RTR^{-1}=T$ and $RJ(n)R^{-1}=J(Rn)$, where $Rn:=RnR^{-1}$.
\item The rotations $\so3_+$ do not commute with $G$; the diagonal rotations do, and are symmetries of the field reading: $\rho_{\rm oct}(L_i)=L_i$
and $\Phi$ is equivariant (\cref{r3:thm:phi}).
\item A rotor $q\in\mathrm{Spin}(V^+)$ with $qe_1q^{-1}=n$ fixes $T$, the $s_a$ and $e_4e_5e_6$, and it carries $J(e_1)$, $\Omega_{e_1}$, $\lint(e_1)$,
$\phi_{e_1}$ to $J(n)$, $\Omega_n$, $\lint(n)$, $\phi_n$.
\end{enumerate}
\end{proposition}

\begin{proof}
$T$ commutes with every same-sign bivector, hence with $t_i$ and $L_i$, and $[Y,Tn]=T[Y,n]$. The spatial part of $L_i$ commutes with $n$, so
$[L_i,n]=-[t_i,n]\in V^+$. A temporal bivector commutes with every $s_a$ and with $e_4e_5e_6$; with $\phi_n(L_i)=s_j\wedge s_k$, $\phi_n(K_i)=s_i\wedge n$ this
gives (iii). (ii) is \cref{r3:lem:rot}. The finite statements are checked with the rational rotors $q=u'(e_1-n)$, $u'\perp n$, at
$n=(2,3,6)/7,(1,2,2)/3,(-1,4,8)/9$ and with an exact Cayley transform of $L_3$ (r4-3, r4-4).
\end{proof}

Fixed. Rotations of the temporal triple move the clock and with it the whole Dirac structure, but they are not symmetries of the field reading;
the diagonal rotations move the clock as well and are symmetries of both.

\begin{proposition}[Permutations of the temporal axes and the reversal of the clock]
\label{r4:prop:perm}
For $\sigma\in S_3$ let $A_\sigma$ be the automorphism of $\Cl(3,3)$ with $e_i\mapsto e_{\sigma(i)}$, $s_a\mapsto s_a$, and $A'_\sigma$ the one with
$e_i\mapsto e_{\sigma(i)}$, $s_a\mapsto s_{\sigma(a)}$; write $\sigma n=\sum_in_ie_{\sigma(i)}$.
\begin{enumerate}[label=(\roman*),leftmargin=2.2em]
\item Both are inner, $A=\Ad_U$ with $U$ unique up to scale and orthogonal up to scale, and
\[
A(T)=\operatorname{sgn}(\sigma)\,T,\qquad A(J(n))=\operatorname{sgn}(\sigma)\,J(\sigma n),\qquad A(n)=\sigma n .
\]
\item $A'_\sigma$ commutes with $G$, permutes the pole operators, $A'(\omega_i)=\omega_{\sigma(i)}$, and with $U|\mathrm{vac}\rangle=|\mathrm{vac}\rangle$ it acts on
the octants by $U\Phi(X)=\Phi(\tilde X)$, $\tilde X=(\operatorname{sgn}(\sigma)P_\sigma\bB;P_\sigma\bE)$, where $P_\sigma$ permutes the components: $\bB$ is
transformed as an axial vector. The transported octant map $\Phi':=U\Phi(P_\sigma\oplus P_\sigma)^{-1}$ is $\Phi$ for even $\sigma$ and
$\Phi\circ(-\eta)=\Phi_{\rm t}=-(-1)^N\Phi$, the textbook-orientation map of \cref{r3:thm:maxwell}(iv), for odd $\sigma$, and its duality operator is
$A'(-T)$: $A'(-T)\Phi'=\Phi'G$.
\item No nontrivial $A_\sigma$ commutes with $G$.
\item Let $R=nu_2$, the rotor of the half-turn about $u_1$. Then $RnR^{-1}=-n$, $RTR^{-1}=T$, $RJ(n)R^{-1}=J(-n)=-J(n)$, and $RT=u_1$: every unit vector
$C$ of the surplus plane, the charge conjugation of Theorem~5.1 of \cite{ES17}, is the half-turn reversing the clock composed with $T$.
\end{enumerate}
\end{proposition}

\begin{proof}
(i) $\Cl(3,3)\cong M_8(\R)$ is central simple (Skolem--Noether); $e_{\sigma(1)}e_{\sigma(2)}e_{\sigma(3)}=\operatorname{sgn}(\sigma)T$ by reordering anticommuting factors,
and $A(J(n))=A(T)A(n)$. (ii) $P_\sigma\oplus P_\sigma$ commutes with the block matrix of $G$. $A'(a^+_k)=a^+_{\sigma(k)}$, so $Uf_k=f_{\sigma(k)}$ and
$Ug_m=\tfrac12\epsilon_{mjk}a^+_{\sigma(j)}a^+_{\sigma(k)}|\mathrm{vac}\rangle=\operatorname{sgn}(\sigma)g_{\sigma(m)}$, since $\epsilon_{\sigma(m)\sigma(j)\sigma(k)}=\operatorname{sgn}(\sigma)\epsilon_{mjk}$.
Hence $\Phi'(Y)=\Phi(\operatorname{sgn}(\sigma)\bB_Y;\bE_Y)$, and $\Phi\circ(-\eta)=-(-1)^N\Phi$ because $(-1)^N$ is $-1$ on the $f$'s and $+1$ on the $g$'s; the duality
relation is the transport of $-T\Phi=\Phi G$. (iii) A temporal permutation alone moves $e_i$ and fixes its partner $e_{i+3}$. (iv) $R$ is an even
element of $\Cl(V^+)$, so it commutes with $T$; $nu_2nu_1u_2=-u_2u_1u_2=u_1$. All of this is checked for the six permutations, alone and paired, and
for three clocks (r4-5, r4-6, r4-7).
\end{proof}

Fixed. The statement asked for is confirmed in the form (i): an odd permutation reverses $T$ and $J(n)$ together, and the paired one, which is
the only one that respects $G$ and the octants, reverses the mass form $T$ and the duality operator $-T$ together. The reason is (ii): an odd
permutation of the axes is an improper rotation, under which $\bB$ is axial; on the octants it changes the orientation convention $\Btext=-\bB$
of \cite{ES17} into the textbook one, for which the duality operator is $+T$ (\cref{r3:thm:maxwell}(iv)). The working note \cite{FMB} defers the
temporal-clock permutations and rotations (its Section~1) and states that reversing $n$ reverses $J(n)$ (its Section~5); the latter is confirmed and
sharpened by (iv): the reversal is a rotation in $\mathrm{Spin}(V^+)$, which leaves $T$, the product of the two reversed objects, fixed.

\begin{table}[t]
\centering
\footnotesize
\setlength{\tabcolsep}{3pt}
\renewcommand{\arraystretch}{1.15}
\begin{tabular}{@{}>{\raggedright\arraybackslash}p{2.6cm}>{\raggedright\arraybackslash}p{1.9cm}>{\raggedright\arraybackslash}p{5.2cm}>{\raggedright\arraybackslash}p{5.0cm}@{}}
\toprule
operation & $T$; $J(n)$ & Dirac reading: rest flow $e^{-m\lambda T}$, senses & field reading: $-T$, $G$, $\Phi$\\
\midrule
rotation $R$ of $V^+$ alone & $T$; $J(Rn)$ & flow invariant; eigenspaces $n=\pm1$ go to $Rn=\pm1$, senses kept & not a symmetry: $[\so3_+,G]\neq0$\\
diagonal rotation & $T$; $J(Rn)$ & as above, with $V^-$ rotated alike & symmetry: $\rho_{\rm oct}(L)=L$, $-T$, $G$, $\Phi$ preserved\\
clock reversal (half-turn) & $T$; $-J(n)$ & flow invariant; the eigenspaces $n=\pm1$ are exchanged, so the particle and antiparticle labels relative to $n$ are exchanged & the octant structure does not involve $n$; the half-turn itself is not a field symmetry\\
odd $\sigma$, temporal alone & $-T$; $-J(\sigma n)$ & $e^{-m\lambda T}\mapsto e^{+m\lambda T}$: sense reversed on each eigenspace; $(m,q)\mapsto(-m,-q)$ & not a symmetry: breaks $G$\\
odd $\sigma$, paired & $-T$; $-J(\sigma n)$ & as for $\sigma$ alone, spatial axes permuted & $G$ kept, octants permuted, $\Phi\mapsto\Phi_{\rm t}$, duality $-T\mapsto+T$\\
$C=RT$, $C\in U_n$ & $T$; $-J(n)$ & same flow; particle $\leftrightarrow$ antiparticle (Proposition~3.8 of \cite{RE}) & not a field symmetry ($C$ is a vector)\\
\bottomrule
\end{tabular}
\caption{What the operations of \cref{r4:prop:rot,r4:prop:perm} do to $T$, $J(n)$ and the two readings (r4-3 to r4-8).}
\label{r4:tab:ops}
\end{table}

\begin{corollary}[The two senses at every clock]
\label{r4:cor:senses}
On the eigenspaces $n=\pm1$ of any clock, $T=\pm J(n)$, so the rest flow $e^{-m\lambda T}$ of Proposition~3.7 of \cite{RE} is $e^{\mp m\lambda J(n)}$, the two
senses of the charge rotation that are particle and antiparticle there (Proposition~3.8 of \cite{RE}). An odd temporal permutation carries the
Dirac data $(K^\mu_n,T,J(n))$, with $K^\mu_n=-n\gamma^\mu_n$, to $(K^\mu_{\sigma n},-T,-J(\sigma n))$: the transported Lagrangian \cref{r1:eq:LD} has mass $-m$ and
charge $-q$. The axial rotation by $\Omega_{\sigma n}$, which preserves every $K^\mu$ and $J$ and reverses $T$ (Proposition~3.6(iv) of \cite{RE}), restores the
mass, leaving $q\mapsto-q$, the effect of $C$.
\end{corollary}

\begin{proof}
With $\Pi_\pm=\tfrac12(\id\pm n)$, $T\Pi_\pm=\tfrac12(T\pm J(n))=\pm J(n)\Pi_\pm$ because $J(n)n=T$. The rest follows from \cref{r4:prop:perm}(i) and from
$\Omega^{\!\top}T\Omega=-T$, $\Omega^{\!\top}K^\mu\Omega=K^\mu$, $[\Omega,J]=0$ (r4-8).
\end{proof}

Proposition~3.7 of \cite{RE} is restated at every clock (fixed); the effect of the permutations on the Dirac reading is fixed; that a Dirac mass
sign is not observable alone is inherited \cite{ItzyksonZuber}. In the decomposition $C=RT$ of \cref{r4:prop:perm}(iv), $T$ preserves the flow,
the eigenspaces of $n$ and $J(n)$ (it is a parity, \cref{r4:thm:roles}), so the reversal of the sense by $C$ in Proposition~3.8 of \cite{RE} comes entirely
from the half-turn that reverses the clock (fixed; the reading of particle and antiparticle as the two clock orientations is interpretive).

\subsection{\texorpdfstring{The three roles of $T$, and $J(n)$ on the octants}{The three roles of T, and J(n) on the octants}}
\label{r4:sec:roles}

\begin{theorem}[The three roles of $T$]
\label{r4:thm:roles}
\begin{enumerate}[label=(\alph*),leftmargin=2.2em]
\item \emph{Mass form.} $\gamma^0_nJ(n)=T$ for every clock: the mass kernel of the Dirac reading is the same at every clock, while the kinetic kernels
$K^\mu_n=TJ(n)\gamma^\mu_n=-n\gamma^\mu_n$ move with it.
\item \emph{Parity.} $\Ad_T|_V=\eta$, the frame parity fixing $\bB$ and reversing $\bE$, and $T$ is up to scale the only element implementing it. At every
clock $T$ commutes with $J(n)$ and $TH_n(\bp)T^{-1}=H_n(-\bp)$ for $H_n(\bp)=\sum_ans_ap_a+mn$: $T=\gamma^0J$ is the admissible Dirac parity with $P^2=-\id$
of Theorem~5.3 of \cite{ES17}, at every clock.
\item \emph{Duality.} $-T\Phi=\Phi G$, $Tf_k=g_k$, $Tg_k=-f_k$, $T|\mathrm{vac}\rangle=|\mathrm{full}\rangle$: on the field octants $-T$ is the Hodge star.
\end{enumerate}
The complex structure $J(n)$ is, for the clock $n$, the selected imaginary unit and the charge generator of the Dirac reading (Theorem~4.1 of
\cite{ES17} at every clock, \cref{r2:thm:i}); in \cref{r1:eq:LD} it enters as $D_\mu=\partial_\mu+qA_\mu J(n)$.
\end{theorem}

\begin{proof}
(a) is \cref{r4:prop:identity} with $\gamma^0_n=n$, and $TJ(n)=-n$. (b) $T$ commutes with $n$ and anticommutes with $s_a$, so it commutes with $J(n)$ and
reverses the momentum term of $H_n$; uniqueness: the solutions of $Ye_a=\eta_{aa}e_aY$ form a line (Skolem--Noether). (c) is
\cref{r3:prop:parity}(i) (r4-9).
\end{proof}

Restated and fixed: (a) is Remark~4.2 and Proposition~4.5 of \cite{ES17} and \cref{r2:thm:i}(iii) at every clock, (b) is Proposition~2.11 of \cite{RE} and
\cref{r3:prop:parity} together with Theorem~5.3 of \cite{ES17} at every clock, (c) is \cref{r3:prop:odd}(i); proved together here (r4-9). The three are
compatible: (b) and (a) say that the mass is parity-even, (b) and (c) that the parity lift acts on the field octants as minus the Hodge star, which is
\cref{r3:prop:parity}. Of the three roles only (c) belongs to the field reading, and it uses $T$ alone; (a) uses $T$ as the product $nJ(n)$.

\begin{proposition}[$J(n)$ relative to the field structure]
\label{r4:prop:Jdec}
\begin{enumerate}[label=(\roman*),leftmargin=2.2em]
\item $J(n)=-L(n)+A(n)$, with $L(n)A(n)=A(n)L(n)=0$; for a clock, $-L(n)^2$ and $-A(n)^2$ are complementary projectors of rank four, $L(n)^3=-L(n)$ and
$A(n)^3=-A(n)$.
\item $L(n)=\rho_n(L(n))-\tfrac12J(n)$, in which $\rho_n(L(n))=\tfrac12\sum_in_is_js_k$ is the Dirac spin rotation about $n$; and $\rho_{\rm oct}(L_i)=\rho_n(L_i)-\tfrac12J(e_i)$ for
every $i$ and every clock. Hence $e^{2\pi L(n)}=e^{2\pi\rho_n(L(n))}e^{-\pi J(n)}=(-\id)(-\id)=\id$.
\item For every $n$,
\[
J(n)\Phi(X)=-\Phi(n\times\bB;\,n\times\bE)+(\bE\cdot n)\,|\mathrm{full}\rangle-(\bB\cdot n)\,|\mathrm{vac}\rangle ,
\]
and $J(n)|\mathrm{vac}\rangle=g(n)$, $J(n)|\mathrm{full}\rangle=-f(n)$. The kernel of $L(n)$ is spanned by the two poles and the longitudinal octants $f(n)$, $g(n)$;
on it $J(n)=A(n)$. On the transverse octants $f(u),g(u)$, $u\perp n$, $J(n)=-L(n)$, the rotation about $n$: $L(n)f(u)=f(n\times u)$, $L(n)g(u)=g(n\times u)$.
\item $J(n)$ does not commute with $\rho_{\rm oct}$ and does not preserve the field octants $\Phi(V)$.
\end{enumerate}
\end{proposition}

\begin{proof}
(i) With $t(n)=\sum_in_ie_je_k=J(n)$ and $s(n)=\sum_in_is_js_k$: $L(n)=\tfrac12(s(n)-t(n))$, $A(n)=\tfrac12(s(n)+t(n))$, $s(n)$ and $t(n)$ commute and
$s(n)^2=t(n)^2=-|n|^2\id$, so $L(n)A(n)=\tfrac14(s^2-t^2)=0$ and $-L(n)^2=\tfrac12(\id+st)$, $-A(n)^2=\tfrac12(\id-st)$ with $(st)^2=\id$. (ii) $\rho_n(L(n))=\tfrac12s(n)$,
and $L_i=\tfrac12s_js_k-\tfrac12e_je_k$ with $e_je_k=J(e_i)$; $\tfrac12s(n)$ and $\tfrac12J(n)$ commute and square to $-\tfrac14\id$. (iii) $J(n)=\sum_in_ie_je_k$ and $e_I|\mathrm{vac}\rangle=a^+_I|\mathrm{vac}\rangle$
(\cref{r3:prop:fock}) give $J(n)|\mathrm{vac}\rangle=g(n)$; $L(n)$ acts on the $f$'s and $g$'s as the vector rotation (\cref{r3:lem:rot}) and kills the poles; $A(n)f(u)=(n\cdot u)|\mathrm{full}\rangle$
and $A(n)g(u)=-(n\cdot u)|\mathrm{vac}\rangle$. (iv) follows from (iii), since $\rho_{\rm oct}$ preserves $\Phi(V)$ (r4-10, r4-11, symbolic in $n$, $\bE$, $\bB$, $u$).
\end{proof}

Fixed. What this means is the following, and the reading is interpretive. The Dirac charge rotation $e^{\theta J(n)}$, read on the octant module, is not
a transformation of the field: on transverse fields it is the rotation by $-\theta$ about the axis paired with the clock, and on the four-dimensional
complement it exchanges the longitudinal components with the pole components, $\bE\cdot n$ with the slot of $\nabla\cdot\bB$ and $\bB\cdot n$ with the slot of
$\nabla\cdot\bE$ (\cref{r3:thm:maxwell}). So the complex structure of the particle reading mixes the field with the Gauss constraints, and it is not in
the commutant of the field reading's Lorentz representation. Item (ii) shows where integer and half-integer spin part on $S$: about the clock axis the
field rotation is the Dirac rotation combined with half a turn of the charge, and the two signs $-1$ cancel. The working note \cite{FMB} asks that
$J(n)$ be analysed separately from $G_{\rm EM}=\star$ (its Section~8, item~5): on the octants $\star$ is $-T$ (\cref{r3:prop:odd}), $J(n)=nT$ differs from it by
the clock, and $J(n)$ commutes with $\lint(n)$ while $G$ commutes with $\lG$.

\subsection{Two Lorentz representations on one module, and their couplings}
\label{r4:sec:reps}

\begin{proposition}[Commutants]
\label{r4:prop:comm}
\begin{enumerate}[label=(\roman*),leftmargin=2.2em]
\item The commutant of $\rho_n(\lG)$, the spin lift of $\lint(n)$, in $\End_\R(S)$ is $\operatorname{span}\{\id,u_1,u_2,J(n)\}\otimes\operatorname{span}\{\id,\Omega_n\}\cong M_2(\C)$,
of dimension $8$.
\item The commutant of $\rho_{\rm oct}(\lG)$ is $\operatorname{span}\{\Pi_F,T\Pi_F\}\oplus\End(\operatorname{span}\{|\mathrm{vac}\rangle,|\mathrm{full}\rangle\})$, of dimension $6$, with $\Pi_F$
the projector onto the field octants.
\item The joint commutant of $\rho_n(\lG)$ and $\rho_{\rm oct}(\lG)$ is $\R\id$, for every clock.
\end{enumerate}
\end{proposition}

\begin{proof}
(i) $u_a$ anticommutes with each vector of $E_n$ and so commutes with its bivectors; $\Omega_n$ commutes with the even elements of $\Cl(E_n)$; the eight
products are independent, and the commutant has dimension $8$ because $S$ complexifies to $2(\tfrac12,0)\oplus2(0,\tfrac12)$, whose commutant is
$M_2(\C)\oplus M_2(\C)$ \cite{GoodmanWallach}. (ii) $\Phi(V)$ is irreducible of complex type with commutant $\operatorname{span}\{\id,T\}$ (\cref{r3:prop:odd}(iii)), the poles are two
trivial summands, and there is no map between them. (iii) Exact computation at $n=e_1,(2,3,6)/7,(1,2,2)/3$ (r4-12); structurally, $\rho_{\rm oct}(L_i)-\rho_n(L_i)=-\tfrac12J(e_i)$
spans $\so3_+$, which with $\lint(n)$ generates $\so{3,3}$, whose commutant $\operatorname{span}\{\id,\Gamma\}$ meets that of the odd $Q_i$ in $\R\id$.
\end{proof}

Fixed; (i) contains the commutant $\HH$ of the four $\gamma^\mu_n$ of \cref{r2:thm:i}(ii). No complex structure, and no operator other than a scalar,
commutes with both Lorentz actions.

\begin{proposition}[Three inequivalent representations on $S$]
\label{r4:prop:ineq}
On the same module $S$ the Lie algebra $\lG\cong\so{1,3}$ acts in three ways: through $\rho_n$ (complexification $2(\tfrac12,0)\oplus2(0,\tfrac12)$, two
Dirac spinors), through $\rho_{\rm oct}$ ($(1,0)\oplus(0,1)\oplus2(0,0)$, the field strength and two singlets) and through the spin lift of $\lG$ itself
($2(\tfrac12,\tfrac12)$). The Casimirs $\sum_k(\rho(K_k)^2-\rho(L_k)^2)$ and $(\sum_k\rho(K_k)\rho(L_k))^2$ are $\tfrac32$ and $-\tfrac9{16}$ for $\rho_n$; $4$ and $-4$ on the field
octants and $0$ on the poles for $\rho_{\rm oct}$; $3$ and $0$ for the spin lift of $\lG$. In $\rho_n$, $\rho_n(L_3)^2=-\tfrac14\id$ and $e^{2\pi\rho_n(L_3)}=-\id$; in the other two
$\rho(L_3)^3=-\rho(L_3)$ and $e^{2\pi\rho(L_3)}=\id$. The three are pairwise inequivalent, and no invertible $U$ satisfies $U\rho_n(Y)U^{-1}=\rho_{\rm oct}(Y)$
for all $Y$.
\end{proposition}

\begin{proof}
Equivalent representations have conjugate Casimirs; the values are computed exactly (r4-13). The decompositions are Proposition~13.12 of
\cite{ES17} and \cref{r3:prop:odd}.
\end{proof}

Restated (the decompositions) and fixed (the Casimirs and the inequivalence). That $\rho_{\rm oct}$ does not commute with the spin lift of
$\lint(e_1)$ (r4-13) is the weakest form of this; the proposition says that the two are not even conjugate. They therefore cannot both be the Lorentz
action of one field on $S$. A theory containing a Dirac field and the field octants must treat them as two fields, $\psi$ in a copy of $S$ with
$\rho_n$ and $\Phi(X)$ in another with $\rho_{\rm oct}$, and the joint action of $Y\in\lint(n)$ is $\rho_n$ on $\psi$ and $\rho_{\rm oct}\circ\phi_n^{-1}$ on the octants.
Since $\Phi$ intertwines $\lG$ on $V$ with $\rho_{\rm oct}$ (\cref{r3:prop:odd}(iii)), a coupling to $\Phi(X)$ is a coupling to $X$ with its own action
$Y\cdot X=[Y,X]$, and conversely.

For $X\in V$ let $\ell(X):=\sum_i(B_iL_i+E_iK_i)\in\lG$, an equivariant isomorphism $V\to\lG$ by the brackets and eq.~(14) of Theorem~12.3 of
\cite{ES17}, and
\[
S_n(X):=\rho_n(\ell(X))=\tfrac12\sum_i\bigl(B_i\,s_js_k+E_i\,s_in\bigr),
\]
the field strength as a bivector of $E_n$; under the dictionary of \cref{r2:prop:dict} it is the Clifford image of $F_X^\sharp=\sum_{\mu<\nu}F_X^{\mu\nu}\gamma_\mu\wedge\gamma_\nu$.

\begin{theorem}[Covariant couplings of the Dirac field to the field octants]
\label{r4:thm:couplings}
Let $\psi$ carry $\rho_n$ and $\Phi(X)$ carry $\rho_{\rm oct}\circ\phi_n^{-1}$.
\begin{enumerate}[label=(\roman*),leftmargin=2.2em]
\item No nonzero bilinear form $\psi^{\!\top}M\Phi(X)$ is invariant, nor one on all of $S$ with $\rho_{\rm oct}$.
\item The invariant trilinear couplings $\psi^{\!\top}M(X)\psi$, $M:V\to\End(S)$ linear with $\rho_n(Y)^{\!\top}M(X)+M(X)\rho_n(Y)+M(Y\cdot X)=0$, are exactly
$M(X)=nC\,S_n(X)$ with $C$ in the commutant of \cref{r4:prop:comm}(i): an eight-dimensional space. For a real Grassmann $\psi$ only antisymmetric
kernels contribute, and these form the two-dimensional space spanned by
\[
M_n(X)=T\,J(n)\,S_n(X)=-n\,S_n(X),\qquad M^\star_n(X)=T\,\Gamma\,S_n(X)=M_n(GX),
\]
and $T\Gamma S_n(X)=TJ(n)\Omega_nS_n(X)$.
For a commuting $\psi$ the symmetric kernels form the six-dimensional space given by
$C\in\operatorname{span}\{u_1,u_2,J(n)\}\otimes\operatorname{span}\{\id,\Omega_n\}$. These dimensions hold at every clock.
\item Explicitly $M_n(X)=\tfrac12\sum_aE_as_a-\tfrac12n\sum_iB_is_js_k$ and $M^\star_n(X)=\tfrac12\sum_aB_as_a+\tfrac12n\sum_iE_is_js_k$. Both commute with $J(n)$ (they carry no
charge); under the parity $\psi\mapsto n\psi$ (or $\psi\mapsto T\psi$), $X\mapsto\eta X$, $M_n$ is even, like the mass form $T$, and $M^\star_n$ is odd.
\end{enumerate}
\end{theorem}

\begin{proof}
(i) The rotation $e^{2\pi\rho_n(L_3)}$ is $-\id$ on $\psi$ and $e^{2\pi\rho_{\rm oct}(L_3)}=\id$ on the octants (\cref{r4:prop:ineq}), so an invariant form satisfies
$\beta(-\psi,\chi)=\beta(\psi,\chi)$; the exact solution spaces are zero at $n=e_1,(2,3,6)/7$ (r4-14). (ii) For every $Y$, $\rho_n(Y)^{\!\top}=-n\rho_n(Y)n$: the rotations
$\tfrac12s_js_k$ are antisymmetric and commute with $n$, the boosts $\tfrac12s_in$ are symmetric and anticommute with it. Hence, for $C$ commuting with $\rho_n$,
\[
\rho_n(Y)^{\!\top}nCS_n(X)+nCS_n(X)\rho_n(Y)=nC\bigl(S_n(X)\rho_n(Y)-\rho_n(Y)S_n(X)\bigr)=-nC\,S_n(Y\cdot X),
\]
using $[\rho_n(Y),S_n(X)]=\rho_n([Y,\ell(X)])=\rho_n(\ell(Y\cdot X))$. The map $C\mapsto nCS_n$ is injective, and the space of invariant maps $V\to S^*\otimes S^*$ has
dimension $8$: complexified, $V=(1,0)\oplus(0,1)$, $\Lambda^2S=(1,0)\oplus(0,1)\oplus6(0,0)\oplus4(\tfrac12,\tfrac12)$ and
$\mathrm{Sym}^2S=3(1,0)\oplus3(0,1)\oplus2(0,0)\oplus4(\tfrac12,\tfrac12)$, so the dimensions are $2$ and $6$ \cite{GoodmanWallach}. The symmetry of $nCS$ follows from the transpose
rule and $C^{\!\top}$: $(nCS)^{\!\top}=-nSnC^{\!\top}n$, which is $-nCS$ for $C=\id,\Omega_n$ and $+nCS$ for $C=u_a,J(n)$ and their products with $\Omega_n$. Finally
$TJ(n)=-n$, $TJ(n)\Omega_n=T\Gamma$, $T\Gamma=-e_4e_5e_6$, and $M_n(GX)=M^\star_n(X)$ by (iii). The exact row reduction of the $2304\times168$ and $2304\times216$ linear
systems at $n=e_1$ and $n=(2,3,6)/7$ gives $2$ and $6$, and the displayed kernels span them (r4-15, r4-16). By \cref{r4:prop:coupclock}(i) the dimensions are
the same at every clock. (iii) $nS_n(X)n$ reverses the electric part; $n\,e_4e_5e_6\,n=-e_4e_5e_6$ (r4-17).
\end{proof}

Fixed: the dimensions, the explicit kernels and their properties. The textbook
comparison is inherited. In complex Dirac theory the Lorentz- and gauge-invariant couplings of $\psi$ to the field strength that are bilinear in
$\psi$ and contain no derivative are the Pauli term $\bar\psi\sigma^{\mu\nu}\psi F_{\mu\nu}$, the coupling of an anomalous magnetic moment introduced by
Pauli \cite{Pauli41}, and the electric-dipole term $i\bar\psi\sigma^{\mu\nu}\gamma_5\psi F_{\mu\nu}$, odd under parity and time reversal
\cite{PeskinSchroeder,ItzyksonZuber}. With $J(n)$ in the place of $i$, $\bar\psi\cdots\psi\leftrightarrow\tfrac12\psi^{\!\top}T\cdots\psi$ as in \cref{r1:eq:LD}, and the real
form $\Omega_n$ of $i\gamma_5$ (Proposition~4.6 of \cite{ES17}), $\sigma^{\mu\nu}F_{\mu\nu}$ becomes a multiple of $J(n)S_n(X)$: $M_n$ is the real form of the Pauli term
and $M^\star_n$ that of the electric-dipole term, up to normalisation and the sign convention for $\gamma^a$; their parities in (iii) agree with the
textbook ones. The Lagrangian density
\begin{equation}
\label{r4:eq:LP}
\mathcal L_{\rm P}=\tfrac12\,\psi^{\!\top}\Bigl(\tfrac{\mu}{c}\,M_n(X_{\rm EM})+d\,M^\star_n(X_{\rm EM})\Bigr)\psi,\qquad X_{\rm EM}=(c\bB;\bE),
\end{equation}
has dimension $\mathrm{J\,m^{-3}}$ in SI units for $[\mu]=\mathrm{J\,T^{-1}}$ and $[d]=\mathrm{C\,m}$, and its two couplings have mass dimension $-1$ in natural units: they
are not renormalisable \cite{PeskinSchroeder}. The arena fixes neither constant (open). The kernel $M_n$ is written through $T$ and $J(n)$ as asked:
$T$ is the adjoint form, $J(n)$ the imaginary unit of $\sigma^{\mu\nu}$, and $S_n(X)$ the field bivector; the dual uses the other $\mathrm{Spin}(3,3)$-invariant
form $T\Gamma$ (\cref{r3:prop:energies}(iii)) and no $J(n)$. With these couplings the Dirac sector reads, at every clock,
\[
\mathcal L_D+\mathcal L_{\rm P}=\tfrac12\,\psi^{\!\top}\,T\,\Bigl[J(n)\Bigl(\gamma^\mu_nD_\mu+\tfrac{\mu}{c}S_n(X_{\rm EM})+d\,\Omega_nS_n(X_{\rm EM})\Bigr)-m\Bigr]\psi,
\]
in the units of \cref{r1:eq:LD}: $T$ is the common left factor of the kinetic, mass and Pauli kernels.

\begin{proposition}[The couplings when the clock moves or the axes are permuted]
\label{r4:prop:coupclock}
\begin{enumerate}[label=(\roman*),leftmargin=2.2em]
\item Let $q\in\mathrm{Spin}(V^+)$ be a rotor with $qe_1q^{-1}=n$. Then, with $X$ fixed,
\[
M_n(X)=qM_{e_1}(X)q^{-1},\qquad M^\star_n(X)=qM^\star_{e_1}(X)q^{-1},\qquad qTq^{-1}=T,\qquad qT\Gamma q^{-1}=T\Gamma :
\]
under $\psi\mapsto q\psi$ the couplings at $e_1$ become those at $n$. The factors $J(n)$, $\Omega_n$, $\gamma^\mu_n$, $S_n$ are conjugated, and $T$, the
clock-independent part of every kernel in the display above, is fixed. The couplings themselves depend on the clock: $M_{e_1}\neq M_{e_2}$.
\item For a diagonal rotation $R$, $M_{Rn}(RX)=RM_n(X)R^{-1}$ and likewise for $M^\star$: clock, spinor and field turn together.
\item For an odd temporal permutation alone, $A_\sigma(M_n(X))=M_{\sigma n}(X)$ and $A_\sigma(M^\star_n(X))=M^\star_{\sigma n}(X)$, while $A_\sigma(T)=-T$. For the paired
permutation, $A'_\sigma(M_n(X))=M_{\sigma n}(\tilde X)$ and $A'_\sigma(M^\star_n(X))=\operatorname{sgn}(\sigma)M^\star_{\sigma n}(\tilde X)$, with $\tilde X$ of \cref{r4:prop:perm}(ii).
\end{enumerate}
\end{proposition}

\begin{proof}
(i) $q$ is orthogonal and conjugates $\rho_{e_1}$ to $\rho_n$ (\cref{r4:prop:rot}(iii)); $S_n$, $J(n)$, $\Omega_n$ follow, and $T$, $\Gamma$ are fixed. (ii) The same with the
spin lift of a diagonal rotation, which also rotates $s_a$ and $X$. (iii) $A(TJ(n))=TJ(\sigma n)$ because the two signs cancel; $A$ fixes $s_a$ and maps $n$ to
$\sigma n$, so $A(S_n(X))=S_{\sigma n}(X)$; $A(\Gamma)=\operatorname{sgn}(\sigma)\Gamma$ for $A_\sigma$ and $\Gamma$ for $A'_\sigma$, and $A'(S_n(X))=S_{\sigma n}(\tilde X)$ because
$s_{\sigma(j)}s_{\sigma(k)}$ is $\operatorname{sgn}(\sigma)$ times the cyclically ordered pair (r4-18, r4-19).
\end{proof}

Fixed. The expectation that the couplings change sign with $T$ under an odd permutation is refuted for the Pauli term and confirmed only for the
dual term under the paired permutation: $M_n=TJ(n)S_n$ contains the product $TJ(n)=-n$, in which the two reversals cancel. What changes sign is the
Pauli term \emph{relative to the mass term}, since $T$ reverses and $M_n$ does not; after the axial rotation that restores the mass
(\cref{r4:cor:senses}), $M_n$ is reversed. Physically the paired odd permutation is an improper rotation, and the electric-dipole coupling is the
parity-odd one (\cref{r4:thm:couplings}(iii)).

\subsection{The two first-order operators and the spatial derivatives}
\label{r4:sec:operators}

The Dirac operator of the clock $n$ is $\gamma^\mu_n\partial_\mu=n\partial_0+\sum_as_a\partial_a$, and $\gamma^0\gamma^\mu\partial_\mu=\partial_0+\sum_a\alpha_a(n)\partial_a$ with
$\alpha_a(n):=ns_a$; the octant Maxwell operator is $\partial_t-D$ with $D=\sum_ae_a\partial_a$ (\cref{r3:thm:maxwell}). In both, $\partial_a=\partial/\partial x^a$ differentiates
along the three spatial coordinates, the lengths on $V^-$, in $\mathrm m^{-1}$.

\begin{proposition}[Two Pauli triples]
\label{r4:prop:triples}
\begin{enumerate}[label=(\roman*),leftmargin=2.2em]
\item $\alpha_1(n),\alpha_2(n),\alpha_3(n)$ are symmetric, square to $\id$ and anticommute; so do $e_1,e_2,e_3$. For any such triple $\pi_a$,
$(\sum_ak_a\pi_a)^2=|k|^2\id$, so both $\partial_0\pm\sum_a\pi_a\partial_a$ factor the wave operator.
\item $J(n)\alpha_a(n)=Ts_a$ and $e_a=-TJ(e_a)$: the Dirac triple is one clock times the spatial generators, $\alpha_a(n)=-TJ(n)s_a$; the octant triple consists of
the three coordinate clocks themselves, $e_a=-TJ(e_a)$, one for each spatial axis.
\item No fixed element $Z$ gives $\alpha_a(n)=Ze_a$ for all $a$, nor $\alpha_a(n)=e_aZ$; in particular none of $T$, $\Gamma$, $\omega_i$, $J(n)$ does.
\end{enumerate}
\end{proposition}

\begin{proof}
(i) $n$ and $s_a$ are orthogonal matrices, $n$ symmetric, $s_a$ antisymmetric, anticommuting: $(ns_a)^2=-n^2s_a^2=\id$ and $ns_ans_b=-s_as_b$. (ii)
$J(n)n=T$, and $TJ(e_a)=T^2e_a$. (iii) $e_a\alpha_a(e_1)$ takes three different values for $a=1,2,3$, and so does $\alpha_a(e_1)e_a$ (r4-20).
\end{proof}

\begin{theorem}[The triples are conjugate only by elements of mixed parity]
\label{r4:thm:conj}
For each clock the $U\in M_8(\R)$ with $U\alpha_a(n)=e_aU$ $(a=1,2,3)$ form an eight-dimensional space, a coset of the commutant $\operatorname{span}\{\id,T\}\otimes\Cl^0(V^-)\cong M_2(\C)$
of the $e_a$; its invertible elements are dense. No nonzero solution is even and none is odd, so no element of the Clifford group, of
$\mathrm{Spin}(3,3)$ or of $\so{3,3}$ relates the two triples. Every invertible solution carries the Dirac representation to the left-multiplication
representation of the temporal Pauli algebra,
\[
U\rho_n(L_i)U^{-1}=-\tfrac12e_je_k=-\tfrac12J(e_i),\qquad U\rho_n(K_i)U^{-1}=-\tfrac12e_i ,
\]
whereas $\rho_{\rm oct}(L_i)=-\tfrac12e_je_k+\tfrac12s_js_k$ and $\rho_{\rm oct}(K_i)=-\tfrac12e_i-\tfrac12s_i\omega$; by \cref{r4:prop:ineq} no $U$ relates $\rho_n$ to $\rho_{\rm oct}$.
\end{theorem}

\begin{proof}
Both triples generate copies of $\Cl(3,0)\cong M_2(\C)$, which has a unique simple module, $\R^4$; $S$ is twice it for each, so intertwiners exist and form
a torsor of the commutant (inherited: Wedderburn's theorem). If $U$ had definite parity, $U\alpha_aU^{-1}$ would be even, but $e_a$ is odd. The exact solution
spaces, their even and odd parts, invertibility of a generic solution and the conjugated generators are computed at $n=e_1$ and $(2,3,6)/7$ (r4-21);
$\rho_n(L_i)=-\tfrac12\alpha_j\alpha_k$ and $\rho_n(K_i)=-\tfrac12\alpha_i$.
\end{proof}

\begin{proposition}[Left and right in the Fock picture]
\label{r4:prop:fockLR}
On $S=\Cl(V^+)|\mathrm{vac}\rangle$ (\cref{r3:prop:fock}) the temporal generators act by left multiplication, $e_a(x|\mathrm{vac}\rangle)=(e_ax)|\mathrm{vac}\rangle$, and the spatial ones by twisted
right multiplication, $s_a(x|\mathrm{vac}\rangle)=(\hat x e_a)|\mathrm{vac}\rangle$, $\hat x=(-1)^{|x|}x$. Hence $\alpha_a(n)(x|\mathrm{vac}\rangle)=(n\hat xe_a)|\mathrm{vac}\rangle$.
\end{proposition}

\begin{proof}
On the eight blades $x=e_I$, $I\subset\{1,2,3\}$ (r4-22).
\end{proof}

Inherited in substance: Chevalley's realisation of $\Cl(V\oplus V^*)$ on $\Lambda^*$ \cite{Chevalley54}, in which one Lagrangian half acts on the left and
the other on the right, and the classical fact that Dirac spinors in the Pauli or space-time algebra transform on one side while the field bivector
transforms on both \cite{Hestenes66,Hestenes67}. The section's statements are fixed in the realisation.

These results answer the remark after \cref{r3:thm:maxwell}. The Dirac reading couples $\partial_a$ to the spatial generators, through $\gamma^0\gamma^a=\alpha_a(n)=ns_a$;
in the Fock picture this is the clock $n$ on the left and the temporal Pauli generator $e_a$ on the right. The octant field reading couples $\partial_a$ to
the temporal generators $e_a$ on the left. In both the index $a$ of the derivative is matched to the index $a$ of a generator by the pairing
$a\leftrightarrow a+3$ of the arena, which is background structure, and the derivative is along the length $x^a$ on $V^-$ in both. The two triples are
conjugate (fixed), by elements of mixed parity only, which are no symmetry of the arena (fixed), and any Pauli triple would make the operator
hyperbolic (fixed). So neither operator shows that spatial differentiation follows from temporal directions: the octant operator uses the three
coordinate clocks $e_a$ as spatial Pauli matrices, one per axis, not one clock, and the statement that space is differentiated ``through time'' is a
reading of the choice of representation (interpretive), not a consequence of either equation; as an inference it is refuted, since conjugating by $U$
rewrites Maxwell's equations with the Dirac triple for the octant map $U^{-1}\Phi$.

\subsection{Varying and permuting the clock}
\label{r4:sec:clock}

\begin{table}[t]
\centering
\footnotesize
\renewcommand{\arraystretch}{1.15}
\begin{tabular}{@{}>{\raggedright\arraybackslash}p{7.3cm}>{\raggedright\arraybackslash}p{7.3cm}@{}}
\toprule
depends on the clock $n$ & independent of the clock\\
\midrule
$\gamma^0_n=n$, $J(n)$, $\Omega_n$; $\lint(n)$, $\phi_n$, the boosts $\rho_n(K_i)$; $\alpha_a(n)$; $K^\mu_n$, $H_n(\bp)$; the Dirac complex structure $\pm J(n)$ and charge;
the eigenspaces $n=\pm1$ that carry the two senses; the couplings $M_n$, $M^\star_n$; the base time $t_n$ &
$T$, $T\Gamma=-e_4e_5e_6=-n\Omega_n$, $\Gamma=J(n)\Omega_n$; the rest flow $e^{-m\lambda T}$; the rotations $\rho_n(L_i)=\tfrac12s_js_k$; $\rho_{\rm oct}$, $\Phi$, $G$, $\lG$, the octants and
pole operators; $D=\sum_ae_a\partial_a$ and the coefficients of the octant Maxwell equation \cref{r3:eq:maxwell}\\
\bottomrule
\end{tabular}
\caption{Clock dependence of the objects of the two readings (r4-2, r4-3, r4-18, r4-23).}
\label{r4:tab:clock}
\end{table}

\begin{proposition}[The field reading does not see the clock; the particle reading does]
\label{r4:prop:clockdep}
The objects of the right column of \cref{r4:tab:clock} are defined without $n$; the octant Maxwell equation at clock $n$ is
$(\partial_{t_n}-D)\Phi(X)=-(\rho|\mathrm{vac}\rangle+\bj\cdot f)$, with the same matrices for every clock, the clock entering only through the time coordinate of the
base $E_n$. The objects of the left column depend on $n$: $n\mapsto J(n)$ is injective, $\lint(n)=\lint(n')$ only for $n'=\pm n$, and $M_{e_1}\neq M_{e_2}$. All of
them are carried from one clock to another by the rotors of $\mathrm{Spin}(V^+)$, which are not symmetries of the field reading, and, together with the field,
by the diagonal rotations, which are.
\end{proposition}

\begin{proof}
By inspection of the definitions in \cref{r3:sec:main}; \cref{r4:prop:identity,r4:prop:rot,r4:prop:coupclock}; $\lint(n)$ determines $E_n$ and hence $\pm n$ (\cref{r2:prop:planes}).
\end{proof}

Fixed. This answers the question of Section~8, item~6, of \cite{FMB}, whether different clocks are physically different sectors, gauge-equivalent
choices or an order parameter: the diagonal rotations map the coupled data at $(n,\psi,X)$ to those at $(Rn,R\psi,RX)$ (\cref{r4:prop:coupclock}(ii)), so for
theories built from the structures of this note a change of clock is a change of frame (fixed as a statement of equivalence; its reading as ``no
sector'' is interpretive); whether a dynamics singles out a clock is open, as in \cref{r2:rem:projection}.

\begin{theorem}[Uniqueness of $T$]
\label{r4:thm:unique}
Let $Y\in\Cl(3,3)$.
\begin{enumerate}[label=(\roman*),leftmargin=2.2em]
\item If $n^{-1}Y$ commutes with $\gamma^0_n,\dots,\gamma^3_n$ for three linearly independent clocks, then $Y\in\R T$. Three clocks of the Fisher--Rao octant suffice,
e.g.\ $e_1$, $(1,2,2)/3$, $(2,3,6)/7$. For two clocks $n,n'$ the solutions are $\operatorname{span}\{T,\,(n\times n')\,e_4e_5e_6\}$.
\item If $n^{-1}Y$ lies, for two independent clocks, in the joint commutant $\operatorname{span}\{\id,\Omega_n,J(n),\Omega_nJ(n)\}$ of $\lint(n)$ and $J(n)$ in which Theorem~4.1 of
\cite{ES17} selects the imaginary unit, then $Y\in\operatorname{span}\{T,e_4e_5e_6\}=\operatorname{span}\{T,T\Gamma\}$, the two $\mathrm{Spin}(3,3)$-invariant forms. Indeed $nJ(n)=T$ and
$n\Omega_n=e_4e_5e_6$ for every clock.
\item Consequently $T$ is, up to scale, the only element of $\Cl(3,3)$ that is the product of the clock with a Dirac complex structure (a square root of $-\id$ in
the commutant of the $\gamma^\mu_n$) for every clock and that does not depend on the clock: the products are $\pm nJ(n)=\pm T$.
\end{enumerate}
\end{theorem}

\begin{proof}
(i) The commutant of the $\gamma^\mu_n$ is $\operatorname{span}\{\id,J(n),u_1\Omega_n,u_2\Omega_n\}$ (\cref{r2:thm:i}(ii)), so $Y\in n\cdot(\ldots)=\operatorname{span}\{n,T,u_1e_4e_5e_6,u_2e_4e_5e_6\}$, using
$nu_a\Omega_n=-u_ae_4e_5e_6$. The four have grades $1,3,4,4$, so for $Y$ fixed the grade-one part lies in $\bigcap\R n=0$, the grade-four part in $(\bigcap U_n)\,e_4e_5e_6$,
which is $0$ for three independent clocks and the line of $n\times n'$ for two. (ii) $n\cdot\operatorname{span}\{\id,\Omega_n,J(n),\Omega_nJ(n)\}=\operatorname{span}\{n,e_4e_5e_6,T,n\Gamma\}$, of grades
$1,3,3,5$; the grades one and five vanish for two independent clocks; $T\Gamma=-e_4e_5e_6$ (r4-1) and $\operatorname{span}\{T,T\Gamma\}$ is the space of invariant forms
(\cref{r3:prop:energies}(iii)). (iii) By Theorem~4.1 of \cite{ES17} at every clock (\cref{r2:thm:i}) the Dirac complex structures are $\pm J(n)$. Exact solutions in r4-23.
\end{proof}

Fixed; the two-clock exception in (i) is a correction to the expectation that two clocks suffice. The second element of (ii), $e_4e_5e_6$, is the
clock times the other complex structure $\Omega_n$, which Theorem~4.1 of \cite{ES17} excludes because it anticommutes with the $\gamma^\mu_n$.

\begin{remark}[The earlier ``covariant mass'' $e_4e_5e_6$]
\label{r4:rem:V}
It was proposed in an earlier, superseded version of the programme that, of the two trivectors anticommuting with the three Dirac matrices
$\alpha_a$, the spatial volume $e_4e_5e_6$ is the scalar mass and $e_1e_2e_3$ its pseudoscalar partner. In the present conventions the first half is
confirmed at every clock: the trivectors anticommuting with $\alpha_1(n),\alpha_2(n),\alpha_3(n)$ are $\operatorname{span}\{T,e_4e_5e_6\}$, the clock-independent space of
\cref{r4:thm:unique}(ii) (r4-24). The assignment is corrected. As the matrix $\beta$ of $H=\sum_a\alpha_ap_a+m\beta$, $e_4e_5e_6=n\Omega_n$ is admissible
(symmetric, square $\id$, commuting with $J(n)$, $H^2=(|\bp|^2+m^2)\id$), but relative to the kinetic kernels with $\gamma^0=n$ it gives the pseudoscalar mass
of Proposition~4.6 of \cite{ES17}, whose Grassmann kernel $J(n)e_4e_5e_6=T\Omega_n$ depends on the clock; the scalar mass has $\beta=n$ and the
clock-independent kernel $J(n)n=T$. The element $e_1e_2e_3$ is not a $\beta$ at all (it is antisymmetric with square $-\id$); it is the kernel.
\end{remark}

Interpretation. In the particle reading the clock and the complex structure both move, and their product $T$ does not; in the field reading $T$ is
the duality on the octants and, as the Chevalley--Mukai pairing, the form pairing $E_k$ with $B_k$ (\cref{r3:prop:energies}). The mass form is the part of
the particle reading that the field reading shares. This is interpretive; what is fixed is \cref{r4:thm:unique} and the fact that $T$ is the only
element common to the structures of both readings that appears in a kernel of the Dirac Lagrangian.

\subsection{Answer}
\label{r4:sec:answer}

\begin{enumerate}[label=(\arabic*),leftmargin=2.2em]
\item \emph{Consistency.} $T=nJ(n)$ holds at every clock and is compatible with all operations considered: $T$ is invariant under all rotations of the
temporal triple and the diagonal rotations, reversed together with $J(n)$, and with the duality operator, by odd permutations, and fixed by the
reversal of the clock (\cref{r4:prop:rot,r4:prop:perm}; fixed). Its three roles, mass form, parity and duality, are compatible (\cref{r4:thm:roles};
restated and fixed). It incorporates consistently into both readings.
\item \emph{The Lorentz representations stay apart.} On $S$ the Dirac representation $\rho_n$ and the octant representation $\rho_{\rm oct}$ are inequivalent,
their joint commutant is $\R\id$, and no operator, $T$ included, relates them (\cref{r4:prop:comm,r4:prop:ineq,r4:thm:conj}; fixed). A single module
carrying both as the Lorentz action of one field is obstructed; the Dirac field and the field octants are two fields.
\item \emph{The bridge as a coupling.} The answer to ``The main obstruction'' of \cref{r3:sec:answer}: $\rho_{\rm oct}$, $\rho_n$ and $\Phi$ coexist as the
actions on two different fields, and the only covariant local couplings without derivatives are trilinear: the Pauli kernel $TJ(n)S_n(X)$ and the
electric-dipole kernel $T\Gamma S_n(X)$, a two-dimensional space at every clock, bilinear couplings being excluded by the $2\pi$ rotation
(\cref{r4:thm:couplings}; fixed, the textbook terms inherited). Since $\Phi$ is an equivalence onto its image, the octants add no coupling that $X$ itself
does not have. $T$ is the clock-independent factor of every Dirac kernel; the couplings are conjugated by the rotation that moves the clock, and an odd
permutation reverses the electric-dipole term with $T$ but not the Pauli term (\cref{r4:prop:coupclock}; fixed, the sign expectation refuted for the
Pauli term).
\item \emph{$J(n)$ on the field.} $J(n)=-L(n)+A(n)$ rotates the transverse octants about the clock axis and maps the longitudinal octants to the Gauss-law poles;
the field rotation about the clock axis is the Dirac rotation combined with half a charge rotation (\cref{r4:prop:Jdec}; fixed; the reading interpretive).
\item \emph{Spatial derivatives.} The remark after \cref{r3:thm:maxwell} is answered by \cref{r4:prop:triples,r4:thm:conj,r4:prop:fockLR}: both operators
differentiate along the lengths of $V^-$, the Dirac reading through the spatial generators and one clock, the octant reading through the three temporal
generators; the triples are conjugate only by mixed-parity elements, and nothing derives spatial differentiation from temporal directions (fixed; the
inference refuted).
\item \emph{The clock.} The field reading is clock-independent and the particle reading is not; $T$ is the unique clock-independent product of the
clock with a Dirac complex structure, and with the second complex structure $\Omega_n$ the space becomes $\operatorname{span}\{T,T\Gamma\}$ (\cref{r4:prop:clockdep,r4:thm:unique};
fixed). Clocks are related by the diagonal rotations, a symmetry of both readings and of the couplings, so a change of clock is a change of frame
(fixed as an equivalence); whether a dynamics selects a clock is open.
\item \emph{What remains for a covariant Dirac--Maxwell theory in the arena.} The minimal coupling through the potential and $J(n)$, $D_\mu=\partial_\mu+qA_\mu J(n)$
(\cref{r1:eq:Ltot,r1:prop:total}; inherited, restated), which uses $A_\mu$ on the base and not the octants; the Pauli and electric-dipole couplings
\cref{r4:eq:LP}, with two constants the arena does not fix (open); the absence of a common Lorentz representation on one $S$ (obstructed at one constituent;
the two readings can be housed together at the mixed second rung of Section~13.6 of \cite{ES17}, at the price stated there, restated); and the
quantisation of either reading (open, as in \cite{FMB}).
\end{enumerate}

\paragraph{On the working note.} Proposition~5.1 of \cite{FMB} is confirmed; its statement that reversing $n$ reverses $J(n)$ is confirmed and
sharpened (\cref{r4:prop:perm}(iv)); its deferred temporal permutations and rotations are carried out (\cref{r4:prop:rot,r4:prop:perm}); its item~5 of
Section~8 is answered by \cref{r4:prop:Jdec} and \cref{r4:sec:reps} ($J(n)$ is not the duality, a common Lorentz representation on $S$ does not exist, and the
cross-sector coupling has the units, covariance and parity stated in \cref{r4:thm:couplings} and \cref{r4:eq:LP}); its item~6 by \cref{r4:prop:clockdep}.
Its requirement that a cross-sector coupling specify symplectic and quantum consistency is not addressed here (open).

{\small\renewcommand{\arraystretch}{1.12}
\begin{longtable}{@{}>{\raggedright\arraybackslash}p{0.09\textwidth}>{\raggedright\arraybackslash}p{0.58\textwidth}>{\raggedright\arraybackslash}p{0.27\textwidth}@{}}
\caption{The claims of this section, with their checks and classes.}\label{r4:tab:claims}\\
\toprule
id & claim & class\\
\midrule
\endfirsthead
\toprule
id & claim & class\\
\midrule
\endhead
r4-1 & conventions: Clifford relations, octants, $T$, $\Gamma$, $T\Gamma=-e_4e_5e_6$ & restated\\
r4-2 & $T=nJ(n)=J(n)n$, $J(n)^2=-\id$ at every clock (\cref{r4:prop:identity}) & restated (\cite{ES17}, \cref{r2:prop:J}); FMB confirmed\\
r4-3, 4 & $T$ invariant, $J(n)$ covariant under $\so3_+$ and $\so3_\Delta$; only $\so3_\Delta$ commutes with $G$ (\cref{r4:prop:rot}) & fixed\\
r4-5, 6 & odd permutations: $T\mapsto-T$, $J(n)\mapsto-J(\sigma n)$; paired: $\Phi\mapsto\Phi_{\rm t}$, duality $-T\mapsto T$ (\cref{r4:prop:perm}) & fixed\\
r4-7 & clock reversal is a half-turn fixing $T$; $C=RT$ (\cref{r4:prop:perm}(iv)) & fixed; FMB sharpened\\
r4-8 & $T=\pm J(n)$ on $n=\pm1$; odd permutation $\Rightarrow(m,q)\mapsto(-m,-q)$ (\cref{r4:cor:senses}) & restated (RE); fixed\\
r4-9 & three roles of $T$; $T$ is the Dirac parity at every clock (\cref{r4:thm:roles}) & restated, fixed\\
r4-10 & $J(n)=-L(n)+A(n)$, $LA=0$; $L(n)=\rho_n(L(n))-\tfrac12J(n)$ (\cref{r4:prop:Jdec}) & fixed\\
r4-11 & $J(n)$ rotates transverse octants, maps longitudinal ones to the poles (\cref{r4:prop:Jdec}) & fixed; reading interpretive\\
r4-12 & commutants $M_2(\C)$, $\C\oplus M_2(\R)$, joint $\R\id$ (\cref{r4:prop:comm}) & fixed\\
r4-13 & three inequivalent Lorentz representations on $S$ (\cref{r4:prop:ineq}) & restated (decompositions), fixed\\
r4-14 & no invariant bilinear coupling $\psi^{\!\top}M\Phi(X)$ (\cref{r4:thm:couplings}(i)) & obstructed\\
r4-15, 16 & trilinear couplings: dimension $2$ (Grassmann), $6$ (commuting); Pauli $TJ(n)S_n$ and dual $T\Gamma S_n$ & fixed; textbook terms inherited\\
r4-17 & Pauli term parity-even, dual parity-odd (\cref{r4:thm:couplings}(iii)) & fixed; inherited in substance\\
r4-18 & couplings conjugated by the clock rotor, $T$ fixed; diagonal covariance (\cref{r4:prop:coupclock}) & fixed\\
r4-19 & odd permutation: Pauli term does not change sign, dual does when paired (\cref{r4:prop:coupclock}(iii)) & fixed; expectation refuted for the Pauli term\\
r4-20 & two Pauli triples, $J(n)\alpha_a=Ts_a$, $e_a=-TJ(e_a)$, no one-sided relation (\cref{r4:prop:triples}) & fixed\\
r4-21 & intertwiners: dimension $8$, none of definite parity; Dirac $\mapsto$ left Pauli representation (\cref{r4:thm:conj}) & fixed; existence inherited\\
r4-22 & temporal $=$ left, spatial $=$ twisted right multiplication (\cref{r4:prop:fockLR}) & inherited (Chevalley), verified\\
r4-23 & uniqueness of $T$; two-clock exception; $\operatorname{span}\{T,T\Gamma\}$ (\cref{r4:thm:unique}) & fixed\\
r4-24 & trivectors anticommuting with $\alpha_a(n)$; $e_4e_5e_6$ as $\beta$ is the pseudoscalar mass (\cref{r4:rem:V}) & fixed; earlier assignment corrected\\
-- & spatial derivatives ``from temporal directions'' (\cref{r4:sec:operators}) & interpretive; inference refuted\\
-- & a change of clock is a change of frame (\cref{r4:prop:clockdep}) & fixed (equivalence); interpretive\\
-- & mass form $=$ the shared part of the two readings & interpretive\\
-- & values of $\mu$, $d$; a dynamics selecting the clock; quantisation & open\\
\bottomrule
\end{longtable}}
